\chapter{The reference implementation}
\label{toolkit:chapter}

The accompanying package turns a precisely specified comparison into a
portable algebraic certificate and a complete experimental plan. Its public
entry point begins with the scientific question, the two target labels, the
available interventions, and the preparation. The program does not infer
these ingredients from a dataset. A scientist supplies them and can inspect
every consequence reported by the calculation.

Version 0.10.0 retains a strict JSON workflow for two finite classical
instruments with rational coefficients. A second, Python-level interface
exposes finite linear representations, continuous-time approximation,
bounded adaptive design, and sequential model checks. These interfaces have
different hypotheses; their availability in one package does not silently
extend the scope of the JSON workflow.

\section{Install and run the first comparison}

From the root of the released reference package, use Python 3.10 or later:
\begin{verbatim}
python -m pip install .
constructive-access analyze examples/passive.json -o passive-report.json
constructive-access analyze examples/active.json -o active-report.json
constructive-access verify examples/active.json active-report.json
\end{verbatim}
The installation requires SymPy and mpmath. The portable algebra checker itself uses
only Python's standard library. For an existing environment with the
dependencies installed, the same commands are available through
\texttt{python -m constructive\_access}. No code imports from an earlier
research directory or requires an unreleased local file.

The passive example has two hidden states, initial distribution
$(1/2,1/2)$, and the same state-revealing read operation under each candidate.
All passive transcript laws coincide. The program returns an exact
equivalence certificate; collecting more passive records does not change
that model conclusion.

The active example also permits an unobserved drive. Its transition matrix
has both rows equal to $(3/4,1/4)$ under candidate 0 and $(1/4,3/4)$ under
candidate 1. Reading costs one operation, driving costs two, and fresh
preparation costs one additional operation per complete record. The
least-action-charge distinguishing experiment is drive followed by read.
It has action charge 3 and total declared charge 4 per record. Its complete
two-outcome laws are
\[
 P_0=(3/4,1/4),\qquad P_1=(1/4,3/4).
\]
The reported total variation is $1/2$ and the one-record equal-prior Bayes
error is $1/4$. The outward-rounded affinity calculation certifies 12
independent complete records as sufficient for average error at most
$1/10$, with declared batch charge 48. This is a sufficient count, not a
minimum sample count. Neither the operation charge nor its minimum is a
measured energy consumption.

The witness contains an action--outcome word, but the implemented experiment
retains \emph{all} outcomes of the associated action sequence. A drive is
not repeated for free until a desired outcome occurs. The report lists the
complete normalized law, making that distinction inspectable.

\section{What a contract contains}

\texttt{load\_contract} accepts a file path, a JSON-compatible mapping, or a
validated contract. It rejects floating-point model coefficients, duplicate
JSON keys, unsupported fields, invalid probability kernels, negative
charges, and mismatched action/outcome alphabets. Integers and exact
rational strings such as \texttt{"1/3"} are accepted. Decimals written as
strings are interpreted as the exact rational number they specify, without
a claim that the physical coefficient is exact.

Each contract contains the following information:
\begin{itemize}
\item A question, target name, and two distinct target labels, ordered with
the candidate models.
\item Two normalized classical instruments: a row initial distribution and
one nonnegative substochastic matrix for each recorded outcome of each
action. Outcome matrices sum to a row-stochastic matrix for each action.
Missing probability mass must be represented by an explicit recorded
outcome, such as abort.
\item A declared additive action charge, units, resource interpretation,
per-record overhead, calibration statement, and included and excluded
items. Energy is an allowed interpretation of a calibrated bound, never
an inference from the word length.
\item A preparation description and an explicit declaration of whether
complete records are independently and identically prepared.
\item A desired equal-prior average error, rational precision settings, and
a cap on explicit transcript enumeration.
\item Optionally, a deterministic finite automaton restricting complete
action sequences. Its transitions cannot depend on postselected outcomes.
\item Optionally, a whole-record total-variation calibration radius for
each candidate, a joint calibration failure probability, and the stated
justification for those inputs.
\end{itemize}

The canonical contract hash binds these assumptions and target meanings to
the report. It identifies the supplied specification. It does not establish
the provenance of a physical calibration, authenticate a laboratory, or
prove that the candidate family contains the true mechanism.

\section{Separate synthesis from checking}

The Python entry points are:
\begin{verbatim}
from constructive_access import load_contract, analyze, verify_report

contract = load_contract("examples/active.json")
report = analyze(contract)
verification = verify_report(contract, report)
assert verification["verified"]
\end{verbatim}
Analysis builds a paired linear representation, runs an exact
charge-priority span search, exports its filtered closure identities, and
checks them. A witness is minimal in the lexicographic order
\((\hbox{additive action charge},\hbox{word length})\).
The report checker independently reconstructs the paired model from the
caller-supplied contract and checks every required expansion identity.
It does not rerun the witness search, compute ranks, or trust the
synthesizer's history. It then reevaluates the complete selected experiment
and its reported statistical quantities. Modified conclusions, unsupported
fields in the report, altered probabilities, missing closure identities,
and certificates for a substituted model fail this check.

There are three distinct mathematical outputs. An equivalence certificate
proves equality for every permitted finite protocol in the declared model.
A witness proves a nonzero probability contrast. A finite-error bound
provides a sufficient number of repeated complete records under its stated
sampling assumptions. A nonzero contrast is not perfect discrimination.

The transcript cap is a computational guard, not an observability theorem.
If explicit enumeration exceeds it, the program retains the exact access
certificate and returns \texttt{transcript\_limit\_exceeded}; it does not
issue a statistical plan. If iid preparation is not declared, the program
reports one-record quantities without a repeated-record guarantee.

\section{Calibration can change the operational answer}

Run the uncertainty example with:
\begin{verbatim}
constructive-access analyze examples/calibrated.json -o calibrated.json
constructive-access analyze examples/calibration_gap_closed.json
\end{verbatim}
The first example assumes whole-record TV radii $1/100$ under both
candidates and a joint calibration failure probability $1/100$. These are
illustrative assumptions, not measurements. For the event
$E=\{y:P_0(y)>P_1(y)\}$, the program evaluates
\[
 a=\max\{0,P_0(E)-r_0\},\qquad
 b=\min\{1,P_1(E)+r_1\},\qquad g=a-b.
\]
When $g>0$, classify using the empirical event frequency with threshold
$(a+b)/2$. Conditional on pre-acquisition calibration information, each
fixed actual hypothesis must produce iid classification records. A
simultaneous calibration event certifies the corresponding conditional
whole-record laws. Classification records cannot be reused to choose or
calibrate the event or threshold.

For a conditional error allocation $\alpha$, the exact implementation
chooses an integer $M$ satisfying
\[
 \frac{Mg^2}{2}\ge L\ge\log(1/\alpha).
\]
Here $L$ is a rational upper enclosure obtained by range reduction and the
positive series
\[
 \log x=2\sum_{k=0}^{n-1}\frac{z^{2k+1}}{2k+1}+R_n,
 \quad z=\frac{x-1}{x+1},\quad
 0\le R_n\le\frac{2z^{2n+1}}{(2n+1)(1-z^2)}
\]
after reducing $x$ to $[1,2]$. Hoeffding's inequality gives a conditional
error at most $\alpha$ for each hypothesis. Adding the stated joint
calibration failure probability gives the unconditional bound. There is no
claim that this conservative event test is optimal.

The supplied example has $g=12/25$, threshold $1/2$, and $M=21$. Its
conditional allocation is $9/100$, and the unconditional bound is
$1/10$. Declared batch charge is 84 operations, including the stated
preparation overhead. The negative control increases both radii to $1/4$;
the robust gap closes and the program returns
\texttt{calibration\_gap\_not\_certified}. This failure concerns the
specified uncertainty certificate. It does not prove that every possible
physical experiment fails.

\section{Advanced interfaces and their boundaries}

The top-level Python package also exposes the following preserved
implementations:
\begin{description}
\item[Finite instruments.] Use \texttt{RationalInstrument} for a normalized
classical instrument. Use \texttt{LinearRepresentation} for a general
weighted representation, and \texttt{ActionDFA} for an action grammar.
The functions \path{compare_models}, \path{minimum_cost_witness}, and
\path{experiment_law} support direct finite-model work. A signed linear
representation is not automatically a physical quantum instrument; a
separate complete-positivity and normalization justification is required.
\item[Continuous time.] \texttt{certified\_wait} encloses a finite CTMC
transition by an exact rational stochastic approximation and a row-TV
error. \texttt{protocol\_tv\_radius} and
\texttt{robust\_word\_gap} propagate the stated approximation error.
Equality of approximate transition laws cannot prove equality of the true
continuous-time laws.
\item[Adaptive design.] \texttt{optimal\_design} optimizes the declared
two-hypothesis classical Bayes risk under a finite integer budget and a
mandatory finite horizon. \texttt{best\_fixed\_design} supplies a matched
open-loop comparison. These routines do not optimize arbitrary quantum
policies or unknown physical models.
\item[Sequential adequacy.] \texttt{AdequacyMonitor},
\texttt{Interval}, and \texttt{interval\_bet} maintain e-process evidence
against declared conditional predictive envelopes. Their validity requires
the appropriate conditional predictions at every step; stationary marginal
intervals cannot be substituted.
\item[Linear targets.] \texttt{identify\_linear\_target} checks the
row-space criterion for a known fingerprint matrix.
\texttt{minimax\_certificate} and
\texttt{verify\_minimax\_certificate} compute and check exact
primal--dual recovery bounds for small rational instances over the full
positive parameter orthant with a coordinatewise deterministic error box.
The generator uses vertex enumeration and is not presented as a scalable
general linear-programming solver.
\item[Counting-channel energy.] \texttt{count\_cap} and
\texttt{verify\_count\_cap} produce and check a rational Chernoff
envelope under a separately supplied predictable compensator bound.
\texttt{declared\_energy\_cap} combines a checked count cap with a
gross-work-per-event bound and an apparatus bound in the same physical
unit. The compensator, work, and apparatus assumptions must be justified
separately from observation-law calibration.
\end{description}
The finite-family experiment-suite workflow is described separately in this
reference. It addresses target-specific design across multiple candidates;
the two-model JSON command remains deliberately explicit about its smaller
input class.

For a concrete linear recovery example:
\begin{verbatim}
from constructive_access import minimax_certificate, verify_minimax_certificate

F = [[1, 1], ["1/2", 1]]
b = [1, 0]
epsilon = ["1/100", "1/100"]
certificate = minimax_certificate(F, b, epsilon)
assert certificate["absolute_error"] == "1/25"
assert verify_minimax_certificate(F, b, epsilon, certificate)["verified"]
\end{verbatim}
This deterministic minimax guarantee assumes a known intervention response
matrix and the specified full error box. It does not certify that a
biochemical reporter implements the declared selective intervention.
For a counting-channel illustration, \texttt{count\_cap(10, "1/100")}
returns 22. With gross event work at most $3/2$ and apparatus work at most
4 in the same unit, the conservative declared energy cap is 37. The example
does not establish those physical input bounds from sensor data.

\section{Reproduction and provenance}

The package contains source-level provenance for every preserved module in
\path{constructive_access/_vendor/PROVENANCE.json}, including the
original SHA-256, the vendored SHA-256, and any import-path substitution.
Source copies are independent of earlier release directories. No software
license is invented or implied by this packaging step.

The workflow tests exercise passive equivalence, active witnesses, charge
minimality checking, zero-charge actions, action grammars, complete
transcript normalization, calibration gap closure, non-iid declarations,
enumeration limits, malformed exact inputs, report forgery, model
substitution, and command-line failure handling. The independent review
tests supplement these with separately derived examples. Run the complete
original workflow test collection using:
\begin{verbatim}
python -m unittest discover -s tests -v
\end{verbatim}
Passing these checks establishes reproducibility for the supported software
contracts. It does not establish a sensor calibration, a successful
laboratory intervention, an unmodeled dissipation estimate, or community
acceptance of the proposed framework.

\section{Response and access-completion modules}
The \path{constructive_access.completion} subpackage adds six modules:
\texttt{physical\_access}, \texttt{quantum\_completion},
\texttt{pulse\_test}, \texttt{graph\_access},
\texttt{robust\_completion}, and \texttt{composition}.
Their contracts and examples appear in the corresponding chapters.
\begin{verbatim}
python examples/completion_tutorial.py \\
    --output results/completion_tutorial.json
python verify.py
python research/clean_reproduction.py
\end{verbatim}
The first command uses the installed package to verify the qubit response
threshold, construct the calibrated three-arm test, and certify both an
energy lower and upper bound of 6400 in the declared model units. The full
reproduction additionally runs producer tests and separate adversarial checks,
regenerates exact certificates and numerical illustrations, and builds this
reference. The clean reproduction builds an offline wheel, installs it in a
separate environment, and checks scientific outputs against the source run.
The software can check a supplied physical contract; it cannot establish an
apparatus calibration or an experimental result by executing a model.

\section{Calibrated explanations and sequential experiments}
The \path{constructive_access.robust} subpackage adds exact physical response
bounds, protected clock design, fixed and anytime statistical tests, and
counts-to-explanation inference. Run the complete installed example with
\begin{verbatim}
python examples/robust_tutorial.py --output results/robust_tutorial.json
\end{verbatim}
The explanation routine receives an already specified preparation, detector,
clock and calibration contract. Its synthetic demonstration is labeled in both
source and output. For actual data the record count must be fixed before
collection, or the separate time-uniform inference method must be used.
The sequential betting fraction and null envelope must be selected before the
current random preparation arm and clock are drawn. Marginal clock agreement
without independence from the physical trajectory is insufficient.

The same example certifies the arbitrary-preparation realization-size witness
and its 6000-record four-arm design. The functions
\texttt{encode\_dimension}, \texttt{dimension\_plan}, and
\texttt{verify\_dimension} preserve the distinction between a complete
witness-error radius and a single arm's error. The quantum certificate uses
exact rational preparations and an independently checked open-system model.


\section{Classical realization and thermodynamic certificates}
The \path{constructive_access.frontier} subpackage exposes five separate
contracts: \texttt{hierarchy}, \texttt{spectral}, \texttt{limits},
\texttt{thermodynamic}, and \texttt{energy\_certificate}. Run their installed
examples together with
\begin{verbatim}
python examples/frontier_tutorial.py --output results/frontier_tutorial.json
\end{verbatim}
The tutorial checks three finite-noise state requirements, an Erlang-clock
certificate, a matching cyclic classical model, a reversible approximation,
a synthetic-count entropy bound, and the continuous-moment curvature witness.
These are different access contracts, not interchangeable black-box tests.

For example, the following inference accepts two predeclared independent-pair
counts. The numbers below are synthetic:
\begin{verbatim}
from fractions import Fraction as F
from constructive_access.frontier import thermodynamic

result = thermodynamic.from_counts(
    5000, 2560, 5000, 3680,
    F(1, 50), F(1, 20), F(8, 5), F(21, 10))
assert result["rate_lower"] > 0
\end{verbatim}
The two lags are $8/5$ and $21/10$ in declared time units, the one-sided
probability radius is $1/50$, and the joint sampling failure allowance is
$1/20$. An entropy interpretation requires stationary sampling, common
homogeneous dynamics, noninvasive fixed readout and ordinary time reversal.
Sampling counts do not validate those physical assumptions.

For a positive-integer lag menu, \texttt{energy\_certificate.certify} accepts
simultaneous probability intervals and positive rational scores. Its exported
certificate records all inputs, a continuous polynomial normalizer bound and
a directed lower bound on relative-entropy separation. The function
\texttt{verify} recalculates that certificate against separately supplied
inputs. It is a deterministic recheck, not a second proof implementation.
Independent tests separately check generator models and polynomial bounds.
The optional \texttt{design} score proposer requires NumPy and SciPy; it is
followed by exact continuous-domain checking. A failed or weak proposal cannot
be interpreted as evidence of equilibrium.

The complete cumulative test and figure collection is run by
\texttt{python verify.py}; \texttt{--skip-pdf} omits typesetting only.
Source provenance binds all five packaged modules to their preserved phase
sources. No previous release directory is required. None of these certificate
functions discovers an apparatus's hidden-state count, reversal parity or
thermodynamic embedding from an undeclared dataset.


\section{Virtual laboratories and apparatus-aware inference}
The \path{constructive_access.virtual} subpackage adds exact contrast and
dependent-record confidence interfaces, timing and exposure controls, and
numerical laboratories for thermodynamic and qubit experiments. Simulations
require the optional numerical dependencies:
\begin{verbatim}
python -m pip install '.[simulation]'
python examples/virtual_tutorial.py --output results/virtual_tutorial.json
\end{verbatim}
The tutorial exercises the installed mathematical interfaces without replaying
the production simulations. The complete frozen simulation contracts are in
\path{virtual/contracts}. Standalone producer scripts and raw retained records
are in \path{virtual/code} and \path{virtual/results}. A full
\texttt{python verify.py} replay regenerates all model experiments, matched
comparisons, exact countermodels, figures and the reference. The separate clean
reproduction installs a wheel and compares scientific outputs byte for byte.

All parameters in these virtual laboratories are illustrative. The simulator knows a hidden
mechanism to evaluate coverage and power; the inference functions do not receive
that mechanism. Exact concentration certificates, numerical physical evolution,
Monte Carlo intervals and measured physical evidence are separately labeled.
Chapter~\ref{ch:ci-empirical} separately reanalyzes published measurements;
those data are included with attribution and an explicit inference limit.

\section{Calibrated inference and public-data reanalysis}
The \path{constructive_access.calibrated} subpackage supplies calibrated
moment supports, nonlinear certificates, activity bounds, dissipation
countermodels, sequential replay, common-block confidence and bounded analog
inference. Run the installed tutorial with
\begin{verbatim}
python examples/calibrated_tutorial.py --output results/calibrated_tutorial.json
\end{verbatim}
Each function retains its own physical and sampling contract. In particular,
stationary marginal records, independent pairs, dependent common blocks and
predictable batches do not share interchangeable confidence corrections.
The nonlinear primal upper bound concerns an observational relaxation, not
an upper bound on physical entropy production. A null upper endpoint in the
resolved-jump entropy helper means positive infinity.

The empirical extra installs HDF5 support. The three licensed, unmodified
measurement files live in \path{calibrated/data/source_audit}; the wheel does
not embed them. Standalone loader and analysis commands reproduce the data
audit and exploratory analysis without network access. Hash and schema
checks precede numerical reads. Neither embedded Python objects nor pickled
force models are evaluated. The empirical chapter records which unavailable
calibrations prevent a thermodynamic certificate.
